\scanpage{38}
en termes de l'invariant $\alpha$ que \deleted{\unsure{cela}} l'isomorphie entre $P$ et $\psi_r(P)$ équivaut à une condition de la forme $\psi_r(\alpha)=\alpha$ [où $\psi_r\in\mathbb Z[\unsure{U|t}]$ est un pol. conv., \deleted{\unsure{\}}} ne dépendant que de $r$, pas de $n$\ldots]

On \deleted{\unsure{trouve|trouvera}} que \deleted{\unsure{l'on a}} si $(n,p)=1$ on aura $\psi_r(\alpha)\neq\alpha$ si $r\neq1$ dans $\Psi_n$, tandis que si \deleted{$\unsure{(n,p)\neq1|\text{complexe}}$} $p\mid n$ (donc ou bien $n=p$, $p$ impair, ou bien $n=2p$) alors \deleted{\unsure{tout|tous}} les $n$-pol. réguliers sur $k$ sont isom. \deleted{Mais} Donc dans ces cas \deleted{\unsure{on a, si}} \inserted{(où on a} $(\mathbb Z/p\mathbb Z)^*=\mathbb Z/(p-1)\mathbb Z$ de card$^{\rm al}$ \inserted{$p-1$} $>2$ \inserted{i.e. $\unsure{\Gamma_n|\Psi_n|r_n}\neq1$} sauf si $p=\unsure{3|2},3$ ( \deleted{le cas $p=2$ étant par ident$^{\rm é}$} \inserted{correspondant aux cas du triangle, du} carré et de l'hexagone) la réponse est négative, sauf dans les seuls \deleted{\unsure{cas}} trois cas envisagés (triangle \inserted{et hexagone} en car. 3, carré en car. 2).

On peut retenir quand $=$ que $\psi_r$ a une nette tendance à changer la classe d'isomorphie d'un $n$-polygône régulier.

On va \inserted{voir si on peut} définir les $\psi_r$ ($r\in\mathbb Z$) \uline{sans restriction sur $r$} pour les repr.
% archive sheet number: 37
